\section{Concentration estimates for the Ricci equation}
\label{sec:parabolic}

We work on the sequence of buffered, parabolically rescaled intervals supplied
by Proposition~\ref{prop:buffered-intervals}.  Write
\[
 I=[-\theta,1+\theta],\qquad \epsilon=\sqrt d\longrightarrow0.
\]
All distances, curvatures, and time variables in this Section refer to these
units.  The number $a\to0$ is constant on each interval and comparable to the
outermost scale at every time.  Sequence indices will usually be suppressed.
Choose a center $x_t$ of the isolated cluster and a nonnegative
$\eta\in C_c^\infty(\operatorname{int}I)$ with $\inf_{[0,1]}\eta>0$.  Set
\begin{equation}
 G=\eta\Ric,\qquad
 K^2=a^{-4}\int_I\int_{B_t(x_t,N_1a)}|G|^2\,d\mu_{g_t}\,dt.
 \label{eq:root-energy}
\end{equation}
The fixed large number $N_1$ will be chosen after the annular parameters below.
All constants may depend on $\eta$.

These integrals can use measurable choices of exact maximizing centers in
Proposition~\ref{prop:root-scale}.  Indeed their nonempty compact sets, in a
fixed closed neighborhood contained in $\mathcal U$, are upper
semicontinuous in time.  To obtain a selection, successively choose the first
member of a finite ordered cover by closed balls of mesh tending to zero that
meets the previously retained compact set, and intersect with that ball.
The resulting compact sets are nested, have diameter tending to zero, and
have a unique common point.  Intersection tests with fixed compact sets are
measurable, so the limiting point is measurable.  All estimates below hold
for every such choice.

\subsection{Geometric inputs and local estimates}
\label{para:inputs}

We specify the consequences of Section~\ref{sec:geometry} that enter the
argument.  They are sequential statements: they apply to arbitrary choices
of times, radii, and windows after passage to a subsequence.
\begin{enumerate}[label=(\roman*)]
\item If $h$ is bounded above and $\epsilon\Delta t/h\to0$, the geometry
at length $h$ moves negligibly on a window of duration $\Delta t$.
On regular charts with $l\ge c h$, its coefficients in fixed material
coordinates change by $O(\epsilon\Delta t/h)$, with any fixed number of
spatial derivatives measured at scale $h$.
\item For $h\asymp a$ on these windows, the old and new centers have
distance $O(a)$.  For $a/h\to0$, they have distance $o(h)$.
For the latter assertion the old cone tip persists as an intrinsic
singularity.  If a concentration representing it were separated from the
current center by a definite fraction of $h$, the limit at their separation
scale would have two distinct genuine exceptional points.  The tree
construction would then give a nonflat ancestor larger than $a$,
contrary to Proposition~\ref{prop:root-scale}.  The cluster itself is fixed
by convergence of the original distances to the isolated point of $Z$.
\item For sufficiently large fixed $C$ and every fixed $C'>0$, distance
annuli with $Ca\le r\le C'$ admit the quotient charts of
Equation~\eqref{geom:uniform-exterior}, with error, through any prescribed finite
order, bounded by a constant times
\begin{equation}
 \Theta(r)=(a/r)^\mu+\delta_i r^\mu,
 \qquad \delta_i\longrightarrow0,
 \qquad \epsilon=o(\delta_i).
 \label{para:theta}
\end{equation}
We decrease $\mu$ when necessary and assume $0<\mu<1$.
The charts cover whole shortened distance annuli, their radial positioning
has arbitrarily small relative error when $C$ is sufficiently large, and
their finite deck groups are nontrivial.
\end{enumerate}
The quantitative estimate in Item (iii) is
Equation~\eqref{geom:uniform-exterior}, with the prescribed finite derivative
order fixed in advance and a fixed positive exterior exponent and constants
uniform in time and radius.  It covers every fixed upper radius $C'$ for
late sequence indices, since $R_i^{\mathrm{out}}\to\infty$.
Whole-band coverage and radial positioning are part of the quotient-chart
construction there.  Positioning is only required with small relative
error, not at the rate $\Theta$.

For a tensor $U$, write
\[
 |U|_{j,r}=\sum_{k=0}^j r^k|\nabla^kU|.
\]
In material coordinates the lower-index Ricci equation and its constraint
are
\begin{align}
 (\partial_t-\Delta_g)\Ric_{ij}
   &=2R_{ikjl}\Ric^{kl}-2\Ric_i{}^k\Ric_{kj},
   \label{para:ricci-equation}\\
 (\partial_t-\Delta_g)G_{ij}
   &=2R_{ikjl}G^{kl}-2\Ric_i{}^kG_{kj}+\eta'\Ric_{ij},
   \label{para:cutoff-equation}\\
 \nabla^i\bigl(G_{ij}-\tfrac12 g_{ij}\tr_gG\bigr)&=0.
   \label{para:bianchi}
\end{align}
The last identity uses that $\eta$ is spatially constant.  On a regular
patch of radius $r$, use time $\tau=(t-t_0)/r^2$ and tensor components whose
size records the original $I$-unit magnitude.  More precisely, for the
spatial chart $\Phi_r$ set
\[
 g_r=r^{-2}\Phi_r^*g,\qquad G_r=r^{-2}\Phi_r^*G,\qquad
 \Ric_r=r^{-2}\Phi_r^*\Ric.
\]
Thus $|G_r|_{g_r}=|G|_g\circ\Phi_r$; the symbol $\Ric_r$ records
the original magnitude and is $r^{-2}\Ric(g_r)$.
We suppress these pullbacks in norm estimates.
Equation~\eqref{eq:ricci-improvement}
gives
\begin{equation}
 |\Ric|_{j,r}\le C_j\epsilon/r,
 \qquad |r^2\eta'\Ric|_{j,r}\le C_j\epsilon r.
 \label{para:source-scale}
\end{equation}
In particular the cutoff source requires no division by $\eta$.

\begin{lemma}[Local smoothing with time $L^p$ input]
\label{para:local-smoothing}
Consider a linear system on a regular parabolic patch, normalized to spatial
radius and temporal duration of order one, with scalar uniformly elliptic
principal part.  Suppose its spatial coefficient derivatives through a
sufficiently high fixed order are bounded.  For nested cylinders
$Q'\Subset Q$ whose evaluation times have a positive past margin, a
solution $U$ with source $F$ satisfies, for $1<p\le2$,
\begin{equation}
 \sup_{Q'}\sum_{k=0}^j|D^kU|
 \le C\left[
 \left(\frac1{|J|}\int_J\|U(\cdot,\tau)\|_{L^2(\Omega)}^p\,d\tau\right)^{1/p}
 +\left(\frac1{|J|}\int_J\|F(\cdot,\tau)\|_{H^{j+3}(\Omega)}^p\,d\tau\right)^{1/p}
 \right],
 \label{para:local-mixed}
\end{equation}
after enlarging $\Omega\times J$ inside $Q$ if necessary.  The first input
can instead be the normalized spacetime $L^p$ norm of $U$.  Constants
depend on the coefficient bounds and margins.  Corresponding time
$L^p$ estimates over longer windows follow by integrating sliding
cylinders of unit duration.
\end{lemma}

\begin{proof}
We give the energy argument, including the change of time exponent.
With a spatial cutoff $\zeta$ and a time cutoff that vanishes on the
past boundary, multiply the equation by $\zeta^2U$ and integrate in
space.  Uniform ellipticity, integration by parts, and Young's inequality
give, on two nested cylinders,
\[
 \sup_\tau\|\zeta U(\tau)\|_2^2
 +\int\|\zeta DU\|_2^2\,d\tau
 \le C_m\int_Q|U|^2
   +C\left(\int_J\|F(\tau)\|_2\,d\tau\right)^2.
\]
Here $C_m$ is bounded by a fixed power of the inverse spatial and temporal
margins.  For the source, pair its time $L^1$ norm with the left-hand
time supremum and absorb half the latter; no time $L^2$ bound on $F$
is used.  First-order and zeroth-order coefficients produce terms
bounded by a small multiple of the gradient energy plus $C\int|U|^2$.

Differentiate spatially.  The top commutator term has the form
$Da\,D^{k+1}U$, paired with $D^kU$, and is absorbed into gradient
energy.  All remaining commutators contain derivatives of $U$ of order
at most $k$ with bounded coefficients.  Applying the preceding estimate
successively on smaller cylinders controls
$\sup_\tau\|U\|_{H^m}$ and $\|U\|_{L^2_\tau H^{m+1}}$ by
$\|U\|_{L^2(Q)}+\|F\|_{L^1_\tau H^m}$.  At each step the spacetime
norm of the next derivative was supplied by the previous gradient
energy.  Taking $m=j+3$ and using the four-dimensional Sobolev embedding
$H^{j+3}\hookrightarrow C^j$ proves the assertion with time exponent
two for $U$ and exponent one for $F$.

For completeness, let $X(\rho)$ denote the resulting spatial and temporal
supremum on a member of a nested family of cylinders.  For
$\rho<\sigma$, interpolate the time $L^2$ spatial $L^2$ norm against
the larger supremum to obtain
\[
 X(\rho)\le C(\sigma-\rho)^{-M}
   X(\sigma)^{1-p/2}\|U\|_{L^p_\tau L^2_x}^{p/2}
   +C(\sigma-\rho)^{-M}\|F\|_{L^p_\tau H^{j+3}_x}.
\]
Normalized and unnormalized time norms are comparable on these fixed
cylinders.  Young's inequality makes the first right-hand term at most
$\kappa X(\sigma)+C_\kappa(\sigma-\rho)^{-M'}\|U\|_{L^p_\tau L^2_x}$.
Take margins proportional to $2^{-n}$ and $\kappa<2^{-M'-1}$.
Iteration sums a convergent geometric series; the terminal term tends
to zero since the smooth solution has finite supremum on the enclosing
compact cylinder.  This proves Equation~\eqref{para:local-mixed}.
Interpolating $\|U\|_{L^2(Q)}$ against the larger spacetime supremum
in exactly this displayed Young inequality, with $\|U\|_{L^p(Q)}$
in place of the mixed norm, proves the spacetime variant.
Finally integrate the $p$th power of the local estimate over the
evaluation time.  Fubini's theorem bounds the integral of the sliding
averages by a constant times the input integral on the enlarged window.
\end{proof}

\subsection{The spherical projection and the missing radial mode}
\label{para:radial-subsection}

Lift a quotient annulus to its Euclidean cover and record tensor
components in a Cartesian frame.  Let $P$ subtract the componentwise
spherical average at every radius.  This operation preserves
equivariance and commutes with the componentwise flat heat operator.

\begin{lemma}[A radial inequality for the oscillation]
\label{para:radial-lemma}
Suppose $PZ=Z$, $Z$ is compactly supported in time, and
$E=(\partial_\tau-\Delta_\delta)Z$.  For $1<p\le2$, set
\[
 z(s)=\left\|\,\|Z(s,\cdot,\tau)\|_{L^2(S^3)}\,\right\|_{L^p_\tau}.
\]
With either ordinary or constant-normalized time measure,
\begin{equation}
 \mathcal Lz:=-z''-\frac3s z'+\frac3{s^2}z
 \le\left\|\,\|E(s,\cdot,\tau)\|_{L^2(S^3)}\,\right\|_{L^p_\tau}.
 \label{para:radial-inequality}
\end{equation}
For $p=2$ the same conclusion holds for a Hilbert direct sum of fields
in different charts or time windows with the same radial variable.
The homogeneous comparison powers are $s$ and $s^{-3}$.
\end{lemma}

\begin{proof}
Put $w(s,\tau)=\|Z(s,\cdot,\tau)\|_2$.  Where $w>0$, pair the
equation with $Z/w$.  Convexity of the Hilbert norm gives
$\langle Z,Z_{ss}\rangle/w\le w_{ss}$, and the zero spherical
mean gives
$\|\nabla_{S^3}Z\|_2^2\ge3\|Z\|_2^2$.  Consequently
\[
 w_\tau-w_{ss}-3s^{-1}w_s+3s^{-2}w\le\|E\|_2.
\]
Multiply by $w^{p-1}$ and integrate in time.  The time term vanishes
by compact support.  If $z=(\int w^p)^{1/p}$, then
\[
 \int w^{p-1}w_s=z^{p-1}z',\qquad
 \int w^{p-1}w_{ss}\le z^{p-1}z''.
\]
For the second inequality, differentiate the first identity and use
$|z'|^2\le z^{2-p}\int w^{p-2}w_s^2$, which is Cauchy's inequality.
H\"older's inequality bounds the integrated source by
$z^{p-1}\|\|E\|_2\|_p$.  Division proves
Equation~\eqref{para:radial-inequality}.  At zeros, work on a finite
time interval containing the support and replace each norm by its
square plus a positive constant under a square root.  The regularized
time boundary values agree, so their derivative integral still
vanishes.  Perform the calculation and let the constant tend to zero
on compact subannuli; the norm and spectral-gap regularization errors
tend to zero there.  The resulting inequalities hold distributionally.
The direct-sum assertion uses the same Hilbert norm computation.
Finally
\[
 \mathcal L(s^\beta)=[3-\beta(\beta+2)]s^{\beta-2},
\]
whose two zero exponents are $1$ and $-3$.
\end{proof}

\begin{lemma}[Control of the spherical mean derivatives]
\label{para:mean-lemma}
Let $T$ be a spatially constant multiple of $\Ric$ on an annulus whose
metric has scaled $C^{j+1}$ distance at most $e$ from $\delta$.
On a fixed-width shell of radius $s$ in its Cartesian coordinates,
\begin{equation}
 \sum_{k=1}^j s^k|\nabla^kT|
 \le C_j\left(\sum_{k=1}^j s^k\sup|D^kPT|
              +e\sup|T|_{j,s}\right),
 \label{para:mean-derivatives}
\end{equation}
where the suprema on the right use a fixed enlargement of the shell.
\end{lemma}

\begin{proof}
Take the flat trace reversal
$S=T-\tfrac12\delta\tr_\delta T$ and write
$S=\overline S(s)+PS$.  Comparing the true divergence constraint
with flat divergence gives
\[
 \partial_iS_{ij}=\mathcal E_j,\qquad
 s^{k+1}|D^k\mathcal E|
 \le C_j e\,\sup|T|_{k+1,s}\quad(0\le k<j).
\]
This follows directly by expanding $g^{ij}-\delta^{ij}$,
$\Gamma(g)$, and the difference between the two trace reversals;
each term contains either a metric error times $DT$ or a metric
derivative times $T$.
For every unit radial vector $n$, the equation at $x=sn$ reads
\[
 \overline S'_{ij}(s)n_i
    =\mathcal E_j(sn)-\partial_i(PS)_{ij}(sn).
\]
Taking $n$ to be each coordinate unit vector bounds every entry
of $\overline S'$.  Radially differentiating this identity at fixed
$n$ bounds the successive derivatives of $\overline S$ by the
displayed error derivatives and derivatives of $PS$.  Cartesian
derivatives of a radial matrix are combinations of these radial
derivatives with powers of $s^{-1}$, so the scaled bounds follow.
In dimension four the trace-reversal map is an involution, since
$\tr_\delta S=-\tr_\delta T$; it is therefore invertible.
Conversion from Cartesian to covariant derivatives adds only the
stated metric-error terms.
\end{proof}

The two Lemmas explain why a degree-zero critical mode cannot obstruct
the estimates.  Although $s^{-2}A$ is componentwise harmonic for a
constant symmetric matrix $A$, its flat trace-reversed divergence is
\[
 -2s^{-3}\bigl(A-\tfrac12\delta\tr_\delta A\bigr)n.
\]
It vanishes for every $n$ only if $A=0$.  Spatially constant tensors
can survive, but have no positive spatial derivatives; their values
will be fixed by an outer boundary estimate.
The trace-reversal calculation is closely related to the improvement
of ALE metric decay in the proof of Deruelle--Kr\"oncke's
Theorem~2.7 in the arXiv version \cite{DeruelleKroncke2021}.
Their field is a stationary
metric perturbation in a chosen gauge. Here the contracted Bianchi
identity constrains the evolving Ricci tensor itself, and the two
lemmas provide the quantitative derivative bounds needed below.

\subsection{The exterior estimate}
\label{para:exterior-subsection}

We apply the missing-mode argument on successive annuli. The goal is
decay faster than the dimensional $r^{-2}$ rate for the part of Ricci
controlled by the root energy. The spatially constant part is instead
controlled from a fixed outer annulus; separating these two contributions
prevents a loss proportional to the number of shells.
Take dyadic radii based on $a$, and define
\[
 \mathcal A_t(r)=\{y:r/2\le d_t(x_t,y)\le2r\},\quad
 A(r)=\left\|\sup_{\mathcal A_t(r)}|G|\right\|_{L^2(I)},\quad
 V(r)=\left\|\sup_{\mathcal A_t(r)}
                   \sum_{j=1}^m r^j|\nabla^jG|\right\|_{L^2(I)}.
\]
Here $m$ is a fixed sufficiently high order, for example $60$.
We may replace widths by fixed neighboring dyadic widths; a star
will indicate a maximum over a fixed number of such neighbors.

\begin{proposition}[Exterior concentration bound]
\label{prop:exterior-ricci}
There are fixed $N_0\gg1$, $r_0>0$, and
$0<\alpha<\min(\mu,1)$, and a choice of $N_1$ in
Equation~\eqref{eq:root-energy}, such that
\begin{align}
 V(r)&\le C\left[(a/r)^{2+\alpha}K
             +\widetilde\delta_i\bigl((a/r)^\alpha+r^\alpha\bigr)\right],
 \label{para:exterior-v}\\
 A(r)&\le C\left[(a/r)^{2+\alpha}K+\widetilde\delta_i\right]
 \label{para:exterior-a}
\end{align}
for $N_0a\le r\le r_0$, where
$\widetilde\delta_i\to0$ and $\widetilde\delta_i\ge\epsilon$.
The constants are independent of $i$ and $r$.
\end{proposition}

\begin{proof}
\emph{Fixed charts and operator errors.}
Fix a large dyadic $H$, to be chosen first, and then a large $T_0$.
Around an intermediate radius $r$ use annuli
$1/(cH)\le s\le cH$ in $r$-units and time windows of length
$T_0r^2$, where $c$ is a fixed ample constant.  Keep each annular
chart fixed in material points.  The motion bound at the smallest
comparable radius is
\[
 C\epsilon T_0Hr\le C_{H,T_0}\Theta(r),
\]
by $\mu<1$ and $\epsilon=o(\delta_i)$.  Thus the metric error through
all required spatial orders is $C_{H,T_0}\Theta(r)$ throughout the
window.  The center and distance comparisons in
Section~\ref{para:inputs} ensure coverage of the same shortened
distance annuli at each time.  Uniform coverage follows by the
sequential criterion, first increasing $N_0$ after $H,T_0$ and then
taking late indices.  Choose $r_0$ small enough that the enlarged
windows near $\supp\eta$ remain in $I$.

Let $\mathcal R(r)$ be the dyadic radii in
$[r/(c'H),c'Hr]$, for another fixed ample $c'$.  Applying
Lemma~\ref{para:local-smoothing} on bounded-width regular patches
and integrating sliding windows bounds true-field derivatives
through any prescribed fixed order by
\begin{equation}
 C_{H,T_0}\left(\max_{\rho\in\mathcal R(r)}A(\rho)+\epsilon r\right)
 \quad\hbox{in }L^2(I).
 \label{para:true-field-bound}
\end{equation}
Regions disjoint from the time support contribute zero.

\emph{Projection and the time cutoff.}
Cover $\supp\eta$ by smaller windows and choose smooth cutoffs
$\chi_\ell$ equal to one there, supported on the respective larger
windows, with uniformly bounded overlap and
$|\partial_t\chi_\ell|\le C/(T_0r^2)$.
Their smaller windows have an additional time margin of order $r^2$.
In the Hilbert direct sum of Lemma~\ref{para:radial-lemma}, use
$Z_\ell=\chi_\ell PG$ in each chart.  Time norms can be measured
in $I$-time: its constant factor relative to $r^2$-time multiplies
both sides of the radial inequality.
Let $\mathcal B_r$ be the dimensionless zeroth-order operator obtained
from the first two terms of Equation~\eqref{para:cutoff-equation}:
\[
 \mathcal B_rU=2\Rm(g_r)(U)-2\Ric(g_r)\circ U.
\]
Its coefficients and scaled spatial derivatives are bounded by
$C_{H,T_0}\Theta(r)$ on the annulus.  With the explicit rescaled
notation above, the flat-equation forcing is the sum of
\[
 \chi_\ell P\bigl[(\Delta_{g_r}-\Delta_\delta)G_r
       +\mathcal B_rG_r+r^2\eta'\Ric_r\bigr]
 \quad\hbox{and}\quad (\partial_\tau\chi_\ell)PG_r,
\]
where the first operator difference includes the connection terms
appropriate to tensor components.
Equation~\eqref{para:true-field-bound} and flatness bound the norm
of the first part, uniformly for $H^{-1}\le s\le H$, by
$C_{H,T_0}(\Theta(r)\max_{\mathcal R(r)}A+\epsilon r)$.
For the second part use the angular Poincar\'e estimate at its
\emph{own} radius:
\[
 \|PG(s)\|_{L^2(S^3)}
 \le Csr\sup|\nabla G|+C_{H,T_0}\Theta(r)\sup|G|.
\]
Bounded overlap is independent of $T_0$.  The cutoff contribution
therefore has the separate bound
\begin{equation}
 \frac{C_H}{T_0}\max_{\rho\in\mathcal R(r)}V(\rho)
   +C_{H,T_0}\Theta(r)\max_{\rho\in\mathcal R(r)}A(\rho).
 \label{para:time-cutoff-term}
\end{equation}
The leading coefficient $C_H/T_0$ contains no chart-persistence
constant.  In particular it does not multiply the uncontrolled
spherical mean of $G$.

At $s=H^{-1}$ and $s=H$ the same angular inequality gives
boundary values at most $C V_*(r/H)$ and $C V_*(Hr)$,
respectively, plus $C_{H,T_0}\Theta\max A$.
The first constant $C$ is independent of $H,T_0$, since each
angular derivative is normalized by the boundary's own radius.
Let $E_0$ denote the sum of the forcing bounds just obtained.
The nonnegative comparison function
\[
 C V_*(r/H)H^{-3}s^{-3}
  +C V_*(Hr)H^{-1}s+\tfrac13H^2E_0
\]
dominates the boundary data after adding the boundary metric errors
to $E_0$, and its $\mathcal L$ is at least $E_0$ on
$[H^{-1},H]$.  The one-dimensional maximum principle, valid also
for the weak inequality, bounds $z$ by this function.  On any fixed
middle annulus this gives the two leading terms with factors
$H^{-3}$ and $H^{-1}$; all inhomogeneous constants may depend on $H$.

Apply flat local smoothing on fixed-width middle patches.  The extra
past margin ensures that each entire unit-size smoothing cylinder lies
where $\chi_\ell\equiv1$.  Thus the equation differentiated in that
cylinder contains only the true operator errors and the $\eta'$ source;
it contains no $\chi_\ell'PG$ term.  Integrating $z^2$ against $s^3\,ds$ on a slightly
larger fixed middle annulus supplies its spacetime input.  Integrate
the resulting local estimates over the smaller time windows and sum
their squares.  The time margins and spatial margins here are of
unit order, so the leading constants are independent of $H,T_0$.
Higher derivatives of the forcing are bounded by
Equation~\eqref{para:true-field-bound} at higher fixed orders.
Finally Lemma~\ref{para:mean-lemma} supplies the positive derivatives
of the spherical mean.  We have proved
\begin{equation}
 \begin{split}
 V(r)\le{}&C\bigl[H^{-3}V_*(r/H)+H^{-1}V_*(Hr)\bigr]
       +\frac{C_H}{T_0}\max_{\rho\in\mathcal R(r)}V(\rho)\\
       &+C_{H,T_0}\left[\Theta(r)
                     \max_{\rho\in\mathcal R(r)}A(\rho)+\epsilon r\right].
 \end{split}
 \label{para:exterior-recurrence}
\end{equation}
It is used only when its enlarged widths stay in the prescribed range.

\emph{Magnitude and boundary estimates.}
At each time join a point on a shell of radius $\rho$ to the next
outgoing shell by a path of length $C\rho$, and join points within
that shell at the same cost.  Connected quotient annuli and whole-band
coverage provide these paths.  Integrating the covariant derivative
of $G$ along the paths, or the derivative of its norm, gives a
pointwise-in-time telescoping bound.  Minkowski's inequality gives
\begin{equation}
 A(r)\le C\sum_{\substack{\rho\text{ dyadic}\\r\le\rho\le r_0}}
                 V(\rho)+C\epsilon.
 \label{para:magnitude-sum}
\end{equation}
A fixed number of neighbors below $r$ may be included to accommodate
the chosen widths.  The outer boundary annuli have $A+V\le
C_{H,r_0}\epsilon$ by Equation~\eqref{para:source-scale}.
The inner boundary annuli have radii comparable to $a$ with fixed,
possibly large, factors.  Choose $N_1$ to include all enlarged
regular patches there at their respective times; this uses the
$O(a)$ center-motion bound.  Local smoothing, followed by integration
of sliding windows, yields
\begin{equation}
 A(r)+V(r)\le C\bigl[(a/r)^2K+\epsilon r\bigr]
 \quad\hbox{on the inner boundary range}.
 \label{para:inner-boundary}
\end{equation}
The factor $r^{-2}$ is the square root of normalized four-dimensional
volume.  Time integration cancels the sliding-window time average,
so it introduces no additional negative power of $r$.

\emph{Closing the recurrence, in the required order.}
Fix $0<\alpha<\min(\mu,1)$, and take
$\widetilde\delta_i\ge\epsilon$ positive and tending to zero
sufficiently slowly.  Put
\[
 W_K(r)=(a/r)^{2+\alpha}K,\quad
 W_0(r)=\widetilde\delta_i\bigl[(a/r)^\alpha+r^\alpha\bigr],\quad
 M=\max_r\frac{V(r)}{W_K(r)+W_0(r)}.
\]
For each index this maximum over finitely many radii is finite.
The geometric series in Equation~\eqref{para:magnitude-sum} imply
\begin{equation}
 A(r)\le C M\bigl[W_K(r)+\widetilde\delta_i\bigr]+C\epsilon.
 \label{para:sum-weight}
\end{equation}
Indeed the three sums involve the powers $-2-\alpha$, $-\alpha$,
and $\alpha$, and are bounded respectively by their initial value,
a constant, and a constant.  Thus no number-of-shells factor occurs.

For each power $r^\beta$ in $W_K+W_0$,
$-2-\alpha\le\beta\le\alpha$.  The two leading terms in
Equation~\eqref{para:exterior-recurrence}, divided by this weight,
are bounded by
\[
 C\bigl(H^{-3-\beta}+H^{-1+\beta}\bigr)M
 \le C' H^{-1+\alpha}M.
\]
Fixed neighboring widths change only $C'$.  Choose $H$ so this
is at most $M/8$.  For this fixed $H$ there is a finite constant
$L_H$ with $\max_{\mathcal R(r)}W/W(r)\le L_H$; choose $T_0$
so $(C_H/T_0)L_H\le1/8$.

For the metric terms in Equation~\eqref{para:sum-weight}, uniform
smallness of $\Theta$ controls the $W_K$ part.  For the constant
$\widetilde\delta_i$ part use the sharper inequality
\begin{equation}
 \frac{\Theta(r)}{(a/r)^\alpha+r^\alpha}
 \le (a/r)^{\mu-\alpha}+\delta_i r^{\mu-\alpha}.
 \label{para:metric-ratio}
\end{equation}
Now increase $N_0$ to make the first term uniformly small on all
needed enlarged radii.  With $r_0$ already chosen for the buffers,
fix $N_1$ large enough for Equation~\eqref{para:inner-boundary}.
Finally take the sequence index sufficiently late to make the second
term small and ensure the stated chart coverage.  This absorbs the
metric contribution into at most $M/4$.  Thus the smallness choices
are made in the order $H$, $T_0$, $N_0$, and the late sequence index;
all constants may depend on the parameters already fixed.
The remaining source is bounded relative to $W_0$ because
$\epsilon r\le\widetilde\delta_i r^\alpha$ for $r<1$.
The two boundary ranges have bounded ratios by
Equation~\eqref{para:inner-boundary} and the outer $O(\epsilon)$
bound.  Their constants may depend on the fixed parameters.
Consequently $M\le C+M/2$.  This proves
Equation~\eqref{para:exterior-v}, and
Equation~\eqref{para:sum-weight} proves Equation~\eqref{para:exterior-a}.
\end{proof}

\subsection{The maximal function controlling the root}
\label{para:maximal-subsection}

Fix a sufficiently large $D_0$ and consider the tree pieces inside
$B_t(x_t,D_0a)$.  Choose an ample fixed $D_1>D_0$, enlarging it to
include the short-time comparison patches.  Extend functions of time
by zero outside $I$, fix $1<p<2$, and set
\begin{equation}
 f(t)=a^{-2}\|G(\cdot,t)\|_{L^2(B_t(x_t,D_1a))},\qquad
 H_*(t)=\bigl(\mathcal M(f^p)(t)\bigr)^{1/p},
 \label{para:maximal-definition}
\end{equation}
where $\mathcal M$ is the centered Hardy--Littlewood maximal operator
on the time line \cite{HardyLittlewood1930}. We include its exact
boundedness argument at the exponent needed below.

\begin{lemma}
\label{para:maximal-lemma}
For every fixed $D_1$,
\begin{equation}
 \|H_*\|_{L^2(\R)}\le C_{D_1,p}(K+o_i(1)).
 \label{eq:maximal-ricci}
\end{equation}
\end{lemma}

\begin{proof}
The energy ball bounds the corresponding part of $\|f\|_2$ by
$K$.  Its complement in the larger ball uses finitely many dyadic
shells, with number and radius factors depending only on $D_1,N_1$.
Volume upper bounds and Proposition~\ref{prop:exterior-ricci} bound
each shell's contribution by
\[
 C(r/a)^2A(r)\le C(a/r)^\alpha K
                       +C(r/a)^2\widetilde\delta_i.
\]
Hence $\|f\|_2\le C_{D_1}(K+o_i(1))$.

We recall the maximal estimate at the exact exponent needed here.
The interval covering argument chooses disjoint witnessing intervals
whose triples cover a compact subset of a maximal level set.  It
gives $|\{\mathcal Mv>\lambda\}|\le C\|v\|_1/\lambda$.
Apply this to $v\mathbf1_{\{v>\lambda/2\}}$, since the discarded
part has every average at most $\lambda/2$.  For $q>1$, integrating
the distribution function then gives
\[
 \begin{split}
 \|\mathcal Mv\|_q^q
 &\le Cq\int_0^\infty\lambda^{q-2}
                     \int_{\{v>\lambda/2\}}v(t)\,dt\,d\lambda\\
 &\le C_q\int_\R v^q.
 \end{split}
\]
Take $v=f^p$ and $q=2/p>1$.  Thus
$\|H_*\|_2\le C_p\|f\|_2$, proving the claim.
\end{proof}

\subsection{A weighted estimate at a single time}
\label{para:interior-subsection}

The exterior estimate controls the collar error in time-integrated
norms. The static input also needs a bound at a single time on every
smaller descendant and joining neck. The maximal function just defined
provides the time control; we now obtain the additional exponential
weight along a deep neck.
\begin{proposition}[Interior estimate with a neck weight]
\label{prop:interior-ricci}
Take any sequence of times with $\eta(t)\ne0$ and extract a tree as in
Theorem~\ref{thm:degeneration}.  On its true compact and half-neck
pieces in $B_t(x_t,D_0a)$, let $r$ be the physical length of the local
regular patch, in $I$-units.  For every fixed required derivative order
$j$ there are $C,\nu>0$ such that
\begin{equation}
 (r/a)^2|G|_{j,r}(t)
       \le C\bigl(H_*(t)+\epsilon\bigr)e^{-\nu D}.
 \label{eq:interior-weighted}
\end{equation}
Here $D=0$ on compact pieces, including truncated outer compact
regions.  On a joining neck $D$, up to bounded shifts, is the minimum
of the logarithmic distances from the two ends.  It is the half-neck
coordinate used in Section~\ref{sec:static-input}.  Constants may depend
on the extracted tree, and the exponent may be decreased.
\end{proposition}

\begin{proof}
For a regular patch of radius $r\lesssim a$, apply
Lemma~\ref{para:local-smoothing} over a past window of length $cr^2$.
The patch lies in the larger ball defining $f$ at its respective
times, by Section~\ref{para:inputs}; shorten its fixed size if needed.
Its normalized spatial $L^2$ input, after multiplying the field by
$(r/a)^2$, is bounded by $Cf$.  The centered maximal average at the
specified time controls its time $L^p$ average.  The source is bounded
by $C\epsilon r(r/a)^2\le C\epsilon$ for late indices.  We obtain
Equation~\eqref{eq:interior-weighted} without its exponential factor.
This proves the compact case and every bounded-depth neck case.

It remains to gain decay as $D\to\infty$.  On such a neck
$r\le Ca e^{-D}$.  Choose $\lambda>0$ sufficiently small and put
\begin{equation}
 h=e^{-4\lambda D},\qquad B=e^{\lambda D}.
 \label{para:neck-window}
\end{equation}
Use the annulus $h\le s\le B$ in $r$-units, with fixed-factor
enlargements, and a time window of length $CB^2r^2$ centered at
the specified time.  By Equation~\eqref{eq:neck-decay}, the initial
spatial error through every required order is bounded by
$Ce^{-\mu(1-4\lambda)D}$.  The worst motion error occurs at the
innermost radius and is at most
\[
 C\epsilon B^2r/h
     \le C\epsilon a e^{-(1-6\lambda)D}.
\]
Choose $\lambda$ so small that
$\min\{\mu(1-4\lambda),1-6\lambda,3-3\lambda\}>3\lambda$.
After decreasing an exponent, the fixed material coordinates
therefore satisfy a uniform scaled error bound
\begin{equation}
 e_D=Ce^{-\delta D},\qquad \delta>3\lambda.
 \label{para:deep-metric-error}
\end{equation}
Regular flat annular limits make the local persistence constants
uniform across these shells.  The entire window stays in $I$,
since $Br\le Ca e^{-(1-\lambda)D}$, and the enlarged spatial patches
remain in the $D_1a$ ball by the same motion comparisons.

Write $T_1=(r/a)^2G$ in components recording unscaled magnitude,
and $Q=H_*(t)+\epsilon$.  For every shell of radius $s$, its
normalized spatial $L^2$ input obeys
\begin{equation}
 (r/a)^2(sr)^{-2}\|G(\cdot,u)\|_{L^2(B_u(x_u,D_1a))}
       =s^{-2}f(u).
 \label{para:interior-normalization}
\end{equation}
Integrate the local estimate on sliding windows of length $c s^2r^2$
inside the larger window.  Fubini's theorem cancels their time
normalization, and the remaining centered window average is bounded
by $H_*(t)$.  Thus, for any fixed high order $J$, both the time
$L^p$ average over a window $J'$ of length comparable to $B^2r^2$ and
the value at time $t$ satisfy
\begin{equation}
 \left(\frac1{|J'|}\int_{J'}\sup_{\text{shell }s}|T_1|_{J,s}^p\,du\right)^{1/p}
      \le CQs^{-2},\qquad
 \sup_{\text{shell }s}|T_1|_{J,s}(t)\le CQs^{-2}.
 \label{para:rough-neck}
\end{equation}
Here derivatives in $|T_1|_{J,s}$ are in $r$-coordinates, so they
represent physical scale $sr$.  To verify the source estimate used
in these applications, its size in shell parabolic units is
\[
 C\epsilon(r/a)^2sr.
\]
After division by $Qs^{-2}$ this is at most
$C(r^3/a^2)s^3\le Ca e^{-(3-3\lambda)D}$.
It is absorbed in Equation~\eqref{para:deep-metric-error}, after
decreasing $\delta$.  This also proves the forcing estimate at
all fixed spatial derivative orders.

Take a time cutoff $\chi$ supported in the large window, equal to
one on an ample middle window, with
$|\partial_\tau\chi|\le CB^{-2}$ in $r^2$-time.  Apply
Lemma~\ref{para:radial-lemma} to $Z=\chi PT_1$, using time measure
normalized by a constant times $B^2$.  The metric and curvature
operator errors cost $CQe^{-\delta D}s^{-4}$: there are two
inverse shell radii from the operator and $Qs^{-2}$ from
Equation~\eqref{para:rough-neck}.  The cutoff costs
$CQB^{-2}s^{-2}$.  Therefore
\[
 \mathcal Lz\le CQ\bigl(e^{-\delta D}s^{-4}+B^{-2}s^{-2}\bigr),\quad
 z(h)\le CQh^{-2},\quad z(B)\le CQB^{-2}.
\]
The four terms in the following barrier have explicit roles:
\begin{equation}
 z(s)\le CQ\left(hs^{-3}+B^{-3}s+B^{-2}
                                  +e^{-\delta D}s^{-2}\right).
 \label{para:interior-barrier}
\end{equation}
The first two are $\mathcal L$-harmonic and dominate the inner and
outer data.  The last two supply the forcing, since
$\mathcal L(B^{-2})=3B^{-2}s^{-2}$ and
$\mathcal L(e^{-\delta D}s^{-2})=3e^{-\delta D}s^{-4}$.
The maximum principle proves Equation~\eqref{para:interior-barrier}.

We next pass from the long time average to the one specified time.
For $1\le s\le B/4$, choose a parabolic patch of temporal size
$cs^2$ where $\chi=1$.  Restricting a normalized $L^p$ time norm
from duration $CB^2$ to this patch costs exactly the bound
\begin{equation}
 C(B/s)^{2/p}.
 \label{para:time-average-loss}
\end{equation}
Angular $L^2$ controls angular $L^p$ because $p\le2$.
Integrating Equation~\eqref{para:interior-barrier} over a fixed-width
radial shell therefore supplies the normalized spacetime $L^p$
input for the flat local estimate in Lemma~\ref{para:local-smoothing}.
The operator errors and their spatial derivatives are controlled
by Equation~\eqref{para:rough-neck}.  Lemma~\ref{para:mean-lemma}
then restores the positive derivatives of the mean, with additional
error $CQe^{-\delta D}s^{-2}$, already bounded by the expression
below.  Since $B^{-3}s\le B^{-2}$ on this range, the result is
\begin{equation}
 \sup_{\text{shell }s}\sum_{k=1}^j s^k|D^kT_1|(t)
 \le CQ(B/s)^{2/p}
        \left(hs^{-3}+B^{-2}+e^{-\delta D}s^{-2}\right).
 \label{para:single-time-derivative}
\end{equation}
Covariant derivatives obey the same bound after the metric-error
terms are included.  The factor in Equation~\eqref{para:time-average-loss}
has been retained in every term.

At radius $B/4$, Equation~\eqref{para:rough-neck} bounds the value
of $T_1(t)$ by $CQB^{-2}$.  Integrate the Cartesian first derivative
along each radial ray down to radius one.  A scaled first derivative
at radius $s$ contributes its bound times $ds/s$; thus
\[
 \begin{split}
 |T_1|_{j,1}(t)
 &\le CQ\left[B^{-2}
   +B^{2/p}\int_1^{B/4}
       \left(hs^{-3-2/p}+B^{-2}s^{-2/p}
                      +e^{-\delta D}s^{-2-2/p}\right)\frac{ds}{s}\right]\\
 &\le CQ\left[B^{-2}+B^{2/p}h+B^{-2+2/p}
                                      +B^{2/p}e^{-\delta D}\right].
 \end{split}
\]
Higher derivatives at radius one are already covered by
Equation~\eqref{para:single-time-derivative}.  The four decay
exponents in $D$ are respectively
\[
 2\lambda,\qquad(4-2/p)\lambda,\qquad
 (2-2/p)\lambda,\qquad\delta-(2/p)\lambda.
\]
They are positive because $1<p<2$ and $\delta>3\lambda$.
Choose $\nu$ smaller than their minimum.  This proves
Equation~\eqref{eq:interior-weighted}.
\end{proof}

There is also an uncut qualitative estimate on these true pieces:
\begin{equation}
 r^2|\Ric|_{j,r}\le C\epsilon r
       \le o_i(1)e^{-\nu D},
 \label{para:uncut-weight}
\end{equation}
with $\nu$ decreased below one, using $r\le Ca e^{-D}$ on deep
necks and $r\le Ca$ on compact pieces.
Equation~\eqref{eq:interior-weighted} supplies the stronger estimate
relative to the energy and cutoff.  In the root units used in
Section~\ref{sec:static-input}, its precise form is
\[
 \eta\,r^2|\Ric|_{j,r}
       \le Ca^2(H_*+\epsilon)e^{-\nu D}.
\]
Together with Equation~\eqref{eq:maximal-ricci}, this is the interface
used for the functional estimate in Section~\ref{sec:contradiction}.
